
\chapter{Revue de littérature}
% Corps du texte
\section{Les troubles du spectre de l'autisme (TSA)}
\subsection{Définition}
Les troubles du spectre de l'autisme (TSA) sont des troubles neurodéveloppementaux caractérisés par des difficultés persistantes dans la communication et les interactions sociales, ainsi que par des comportements et intérêts restreints ou répétitifs \parencite[p.~56–58]{apa_diagnostic_2022}.
Ce trouble apparait pendant la période du développement, en général à la petite enfance. L'appellation "spectre" reflète une forte hétérogénéité individuelle au-delà des critères centraux de la communication et des comportements répétitifs \parencite{oms_classification_2018}. Ainsi, les TSA peuvent également se manifester à travers une déficience intellectuelle, des troubles du langage ou des troubles moteurs \parencite[p.~62]{apa_diagnostic_2022}, et sont souvent accompagnés de comorbidités \parencites[p.~67–68]{apa_diagnostic_2022}{lord_autism_2018}.

%\subsection{Etiologie}
%Rajouter rapidement Etiologie ??

%Principalement génétique, héréditaire à 80\,\%
%+ facteurs environnementaux

\subsection{Prévalence et comorbidités} \label{sec:prev_comor}
La prévalence des TSA est estimée à environ 1\,\% dans le monde, bien que les chiffres varient selon les études et aient tendance à augmenter avec le temps en raison d'une hausse des demandes de diagnostic \parencite{zeidan_global_2022, brugha_epidemiology_2011, Elsabbagh2012Apr}.
Les garçons sont généralement plus nombreux que les filles, avec un ratio masculin/féminin variant autour de 4 \parencite{zeidan_global_2022}, malgré des biais genrés compliquant l'évaluation (voir Section~\ref{sec:norme}). Les TSA seraient associés à une probabilité accrue de déficience intellectuelle (DI), mais la proportion n'est pas connue \parencite{zeidan_global_2022}.
%est importante puisque les chiffres vont de 0\,\% à 70\,\% \parencite{zeidan_global_2022} + 2e source.

Par ailleurs, les TSA s'accompagnent fréquemment de comorbidités psychiatriques. Chez l'enfant, environ 45\,\% présentent par exemple au cours de leur vie un trouble déficitaire de l'attention avec hyperactivité (TDAH), 42\,\% un trouble anxieux et 14\,\% une dépression \parencite{Micai2023Dec}. Chez l'adulte, ces proportions sont respectivement de l'ordre de 22\,\%, 28 à 42\,\% et 34\,\% \parencite{hollocks_anxiety_2019, Micai2023Dec} et sont donc supérieures à celles de la population générale. On y retrouve en effet autour de 4\,\% pour le TDAH \parencite{Popit2024Oct, Song2021Feb}, 9\,\% pour le trouble anxieux \parencite{Kessler2012Aug} et 29\,\% pour la dépression \parencite{Kessler2012Aug, Hu2022May}.

\medskip

La prévalence et la diversité des comorbidités psychiatriques observées invitent donc à un diagnostic le plus précoce possible. En ce sens, quand est-ce que les premières caractéristiques autistiques émergent-elles ? Demeurent-elles stables tout au long du développement ?

\subsection{Apparition du phénotype et trajectoires développementales}
Dès les premiers mois de vie, des signes précoces du phénotype autistique apparaissent. Alors, dès l'âge d'un an, peuvent être observées des difficultés dans les premières interactions sociales, ainsi que des comportements répétitifs ou stéréotypés \parencite{Webb2009Apr,Tanner2021Mar,Zwaigenbaum2015Oct}. Un contrôle moteur atypique, que ce soit au niveau de la motricité globale ou de la  motricité fine, est également décrit dès 6 à 12 mois \parencite{Estes2015Jul,Zwaigenbaum2015Oct}. 

Pour autant, les trajectoires développementales sont hétérogènes. En effet, si le diagnostic de l'autisme est stable dans le temps, on observe une grande hétérogénéité entre les études \parencite{Bieleninik2017Sep}. De plus, bien que l'expression du phénotype autistique tende, en moyenne, à diminuer de l'enfance à l'adolescence \parencite{Waizbard-Bartov2022Nov}, entre 20\,\% et 49\,\% des enfants avec un TSA présentent néanmoins une régression dans les domaines du langage ou de la communication sociale \parencite{Webb2009Apr}. Cette évolution dépend alors de variables comme l'environnement ou la présence ou non de déficience intellectuelle \parencite{Waizbard-Bartov2022Nov}.

\medskip

Ainsi, l'apparition précoce du phénotype autistique et son évolution avec l'âge invitent à s'intéresser à la population pédiatrique. Par ailleurs, la présence d'un éventail de caractéristiques aussi diversifiées dès les premiers stades soulève la question de leurs bases neurobiologiques. Quelles altérations cérébrales pourraient expliquer la diversité des manifestations observées chez les enfants avec TSA ?

\section{Bases neuronales des TSA} \label{sec:bases_neuro}

Les bases neuronales des TSA peuvent être comprises non pas comme l'altération d'une région cérébrale unique, mais comme un ensemble d'anomalies fonctionnelles et morphologiques affectant le cerveau à une échelle globale. Sans prétendre à l'exhaustivité, nous présenterons ici un aperçu des connaissances actuelles afin de situer les altérations cérébrales susceptibles de sous-tendre les caractéristiques observées chez les enfants avec TSA.

\subsection{Bases structurelles} \label{sec:bases_struct}

Sur le plan structurel, les analyses de morphométrie voxel‐à‐voxel (VBM) mettent en évidence un ensemble d'altérations de la matière grise chez les personnes présentant un TSA. Il a notamment été rapporté des réductions de volume dans le cortex cingulaire antérieur et médio-préfrontal, dans l'insula ainsi que dans le cervelet postérieur \parencite{Guo2024Apr}. Ces anomalies sont souvent convergentes entre mesures fonctionnelles et morphologiques \parencite{Patriquin2016Oct}. De plus, il a été mis en évidence des volumes plus faibles du putamen, de l'amygdale, du noyau accumbens et du pallidum. Ces différences sous-corticales semblent relativement stables au cours du développement, tandis que les altérations corticales tendent à s'atténuer à l'âge adulte \parencite{vanRooij2017Nov}.

\subsection{Altérations fonctionnelles}

En imagerie fonctionnelle, les études au repos montrent, chez les sujets autistes, notamment une tendance à une hypoactivation dans l'insula gauche, du cortex cingulaire antérieur et médial préfrontal (ACC/mPFC), ainsi que dans l'aire temporale inférieure droite, accompagnée d'une hyperactivation du précuneus et de l'aire motrice supplémentaire \parencite{Guo2024Apr}. Des altérations concernent également le \textit{default mode network} (DMN), où la connectivité entre le cortex préfrontal et le cortex cingulaire postérieur/precuneus apparaît réduite chez les individus avec TSA, ce qui pourrait expliquer certains troubles dans la cognition sociale \parencite{Sato2019Aug}. Des anomalies comparables touchent aussi les régions du "cerveau social" hors DMN, notamment l'amygdale \parencite{Sato2019Aug}, dont la connectivité affaiblie avec le cortex préfrontal ventromédian est corrélée à la sévérité des symptômes socio-affectifs \parencite{Kleinhans2015Dec}. On observe également une dépendance du profil d'activation en fonction du type de tâche mobilisée : les tâches non sociales montrent un recrutement anormal de l'ACC dorsal et de l'aire motrice supplémentaire, tandis que les tâches sociales révèlent par exemple une hypoactivation du cortex cingulaire postérieur 
\parencite{DiMartino2008Nov}.

Le chevauchement des différences neuronales avec des régions impliquées dans le contrôle moteur, telles que le cervelet, les ganglions de la base et le cortex préfrontal, pourrait expliquer les difficultés motrices fréquemment observées chez les sujets avec TSA. En effet, des lésions du cervelet pourraient perturber le contrôle en  \textit{feedforward} ou la coordination sensori-motrice \parencite{Manto2012Jun}. De même, les ganglions de la base pourraient être liés à l'apparition de mouvements stéréotypés \parencite{Groenewegen2003}. %lié à tourettes's

\medskip
Ainsi, les données de neuroimagerie révèlent des caractéristiques neuronales distribuées plutôt que localisées, affectant des régions impliquées dans la cognition sociale, la régulation émotionnelle et le contrôle moteur. Le chevauchement de ces anomalies avec celles observées dans d'autres troubles du neurodéveloppement ou pathologies complexifie leur interprétation \parencite{Guo2024Apr, Lukito2020Mar}.
Parmi les régions altérées, on retrouve de façon récurrente des structures en lien avec la motricité, notamment le cervelet ou les ganglions de la base. Cette implication dans un trouble défini par des critères socio-cognitifs interroge la place des troubles moteurs dans le phénotype autistique. Alors, comment les connaissances issues des modèles cérébraux actuels s'intègrent-elles, ou non, dans les outils de diagnostic existants, et quelles en sont les limites ?

\section{Diagnostic des TSA : outils, validité, et limites méthodologiques}
Malgré la mise en évidence progressive de différences neuronales suggérant l'existence de marqueurs non seulement sociaux mais aussi moteurs chez les personnes présentant un TSA, la recherche et la pratique clinique se sont historiquement centrées sur les dimensions psycho-socio-cognitives du trouble \parencite{cook_movement_2016}. Pour autant, le diagnostic des TSA reste complexe en raison d'un chevauchement symptomatique (voir section~\ref{sec:prev_comor}) et neuronal (voir section~\ref{sec:bases_struct}) avec d'autres troubles du neurodéveloppement ou d'autres pathologies. 

Actuellement, le diagnostic repose sur une évaluation clinique intégrant l'observation comportementale, l'histoire développementale et les informations parentales, selon un consensus pluridisciplinaire considéré comme le \textit{gold standard} \parencites[p.~52]{HAS2018_TSA}{Randall2018Jul, Hirota2023Jan, falkmer_diagnostic_2013, Woolfenden2012Jan}. Un nombre important d'outils standardisés pour assister le diagnostic ont été proposés \parencites[p.~51]{HAS2018_TSA}{falkmer_diagnostic_2013, Randall2018Jul}. Ainsi, aujourd'hui, les outils standardisés considérés comme la référence et les plus utilisés sont l'\textit{Autism Diagnostic Interview, Revised} (ADI-R) et l'\textit{Autism Diagnostic Observation Schedule, 2e édition} (ADOS-2) \parencite{Hirota2023Jan, frigaux_ladi-r_2019}.



\subsection{Outils comportementaux de référence} %(ADOS/ADI-R)}

L'ADOS-2 est un outil standardisé et semi-structuré d'observation directe, basé sur l'évaluation de la communication, l'interaction sociale, le jeu et les comportements répétitifs chez les personnes suspectées de TSA \parencite{Lord2012ADOS2}. Il se compose de 4 modules adaptés à l'âge et au niveau de langage de l'individu \parencite{Lord2012ADOS2}. L'ADI-R, quant à lui, prend la forme d'un entretien semi-structuré mené auprès des parents ou aidants d'enfants et d'adultes, portant sur la trajectoire développementale de l'individu \parencite{Lord1994Oct}. 

\medskip

Il convient alors de s'interroger sur la validité de ces outils diagnostiques : le statut de référence qui leur est attribué est-il justifié par leur validité ?

\subsection{Validité et questionnement du statut de référence}

Chez les enfants, l'ADI-R possède une sensibilité de l'ordre de 60 à 70\,\% et une spécificité d'environ 80\,\% \parencite{Lebersfeld2021Nov, Fulceri2025Jun, dosSantos2024Screening, Randall2018Jul}. Cependant, ces résultats cachent une incertitude très élevée, avec des plages de sensibilité s'étendant de 19 à 100\,\% et une spécificité comprise entre 50\,\% et 100\,\% \parencite{Lebersfeld2021Nov, Fulceri2025Jun}. L'ADOS-2 montre dans la même population une sensibilité légèrement plus importante, de l'ordre de 90\,\%, et une spécificité autour de 70-80\,\% avec incertitude notable mais plus faible que pour l'ADI-R (sensibilité : 77–100\,\% ; spécificité : 44–100\,\%) \parencite{Lebersfeld2021Nov, Fulceri2025Jun}. %L'usage combiné des deux outils pourrait améliorer la validité au prix d'une baisse de sensibilité. C'est en effet ce qui est observé pour la combinaison de l'ADI-R et de l'ADOS (1ère édition) \parencite{Randall2018Jul, Lebersfeld2021Nov}.

En termes de fiabilité inter‑juges, des accords élevés sont rapportés pour l'ADI-R, avec des taux compris entre 80 et 90\,\% \parencite{Lord1994Oct, Zander2017}. Il en va de même pour l'ADOS-2, bien que la fiabilité dépende fortement de l'expérience de l'évaluateur, avec des taux d'accord variant de 50 à 86\,\% selon les modules et l'expérience de l'évaluateur \parencite{Kamp-Becker2018Sep}.

Un biais circulaire doit également être pris en compte : les études de validation de ces outils s'appuient généralement sur un diagnostic de référence posé par une équipe pluridisciplinaire, laquelle intègre déjà les résultats de l'ADOS/ADOS-2 et de l'ADI-R dans son évaluation.

Enfin, les performances varient aussi selon les profils cliniques : la spécificité de l'ADI-R peut, par exemple, tomber à 57\,\% chez les adultes sans déficience intellectuelle \parencite{adamou_predicting_2021}. Ainsi, bien qu'efficaces, les outils standardisés sont sujets à une validité incertaine et hétérogène en fonction des populations. Cette hétérogénéité questionne la dépendance de ces outils à des normes développementales, mais aussi sociales, qu'il convient d'interroger.

\subsection{Ancrage normatif des outils de diagnostic}\label{sec:norme}

Les outils de diagnostic possèdent un biais de genre marqué : les filles sont régulièrement sous-diagnostiquées \parencite{Gupta2025Jan, Lai2023Dec, Rea2023Jul}. Ce phénomène s'explique en partie par des facteurs socioculturels et par la construction genrée des outils et des critères diagnostiques. En effet, la plupart des tests ont été élaborés à partir de cohortes majoritairement masculines, et certains comportements typiques des TSA peuvent être socialement tolérés, voire valorisés, chez les filles \parencite{Lai2023Dec}. De plus, la difficulté d'application des outils standardisés auprès de populations présentant des particularités sensorielles, telles qu'une déficience visuelle ou auditive \parencite{Hirota2023Jan}, ou encore la nécessité de traduire ces outils et d'en valider ces traductions constituent des limites supplémentaires.

\medskip

En outre, les outils actuellement utilisés pour le diagnostic des TSA présentent des biais %méthodologiques et sociaux 
qui conduisent à une sous-identification récurrente de certaines populations, notamment féminines, et une validité incertaine. Dans cette perspective, une approche fondée sur le contrôle moteur, en complément des outils actuellement utilisés, pourrait contribuer à réduire certains de ces biais. En effet, la motricité et notamment la motricité fine est moins susceptible d'être influencée par les différences de genre, de handicap, ou de contexte culturel que les interactions sociales. De plus, les altérations motrices observées chez les personnes autistes sont corrélées avec les caractéristiques centrales des TSA comme les difficultés sociales ou de communication \parencite{hocking_what_2017}. 

\section{TSA et contrôle moteur : caractérisation de la motricité fine}
\label{sec:controle_moteur}
Dès les premières étapes du développement, les enfants présentant un TSA montrent majoritairement des déficiences dans les évaluations standardisées du développement moteur, qu'il s'agisse de la motricité globale ou de la motricité fine \parencite{hocking_what_2017}. %Des particularités motrices ont notamment été rapportées dans des tâches mobilisant la motricité fine, telles que l'écriture ou le pointage et la saisie manuelle \parencite{hocking_what_2017}.
Toutefois, tous les enfants avec TSA ne présentent pas nécessairement de déficiences motrices \parencite{hocking_what_2017}. Ainsi, plutôt que de ne prendre en compte uniquement ces difficultés motrices, il apparaît pertinent d'examiner les particularités de la motricité globale comme fine chez les enfants avec TSA. Dans cette optique, nous faisons ici le choix de nous intéresser à la motricité fine, à travers l'étude des modèles de pointage, de saisie manuelle et d'écriture.
%il apparaît plus judicieux dans une perspective diagnostique et de compréhension d'évoquer les caractéristiques générales et atypiques de la motricité fines des TSA. Nous aborderons alors modèles de pointage/écriture ?

\subsection{Pointage et saisie manuelle}\label{sec:pointage}
%def pointage ?

La littérature met en évidence, chez les enfants présentant un TSA, une tendance à une augmentation du temps de mouvement lors des tâches de pointage et de saisie manuelle \parencite{benchekri_mieux_2024, Glazebrook2006Jul, forti_motor_2011}, bien que ce résultat ne soit pas systématiquement observé \parencite{Papadopoulos2012Jan, Rinehart2006Aug}, ainsi qu'une diminution de la fluidité du mouvement \parencite{benchekri_mieux_2024}. Alors, cette diminution de la vitesse pourrait être interprétée soit comme la cause d'une réduction de la fluidité du mouvement, soit, au contraire, comme sa conséquence, à l'image de ce qui a été observé chez des adultes atteints du syndrome parkinsonien \parencite{benchekri_mieux_2024}.
De plus, il a été rapporté une augmentation du temps d'élaboration de la commande motrice \parencite{Glazebrook2006Jul, Rinehart2006Aug}, pouvant refléter un dysfonctionnement des ganglions de la base similaire à celui observé chez les patients parkinsoniens \parencite{Rinehart2006Aug}.

Il est important de noter que ces différences sont modulées par des variables de confusion, notamment le quotient intellectuel (QI) ou la complexité de la tâche \parencite{Mari2003Feb}. Ainsi, les enfants présentant un TSA associé à un QI plus faible que la moyenne tendent à exécuter les mouvements plus lentement que ceux ayant un QI plus élevé \parencite{Mari2003Feb}.


\subsection{Écriture} \label{sec:ecriture}
%def écriture, stroke ?
Chez les enfants présentant un TSA, on peut retrouver des difficultés marquées dans l'écriture. On retrouve notamment une diminution de la lisibilité globale et une détérioration de la formation des lettres \parencite{Kushki2011Dec, Handle2022Oct}. Au niveau moteur, l'écriture se caractérise par une vitesse d'exécution ralentie et une fluidité de mouvement plus faible \parencite{benchekri_mieux_2024, kangarani-farahani_motor_2024, Handle2022Oct}.
Il peut être considéré que ces caractéristiques résulteraient notamment d'une altération de la coordination visuomotrice, les enfants présentant un TSA s'appuieraient en effet préférentiellement sur le \textit{feedback} proprioceptif \parencite{benchekri_mieux_2024, kangarani-farahani_motor_2024, shafer_visual_2021}. %L'augmentation du contrôle en \textit{feedback} entraînant plus de corrections et une donc une diminution de la fluidité.
\medskip


Nous avons pu voir que des caractéristiques spécifiques de la motricité fine chez les enfants autistes sont mises en évidence au sein de tâches de pointage ou d'écriture, et ne se limitent pas aux déficits fonctionnels observés par les tests standardisés. Une approche consistant d'abord en une classification, puis en la détermination des profils moteurs, pourrait permettre une caractérisation plus fine de ces différences. Cette démarche contribuerait à la compréhension du phénotype moteur des TSA ainsi que des mécanismes neurofonctionnels sous-jacents. En effet, ces caractéristiques motrices, cohérentes avec les altérations des structures impliquées dans le contrôle moteur (voir section \ref{sec:bases_neuro}), apparaissent comme des manifestations potentiellement informatives du fonctionnement cérébral des TSA. Une description précise de ces profils moteurs pourrait ainsi favoriser l'identification de marqueurs pertinents, utiles tant pour améliorer le diagnostic précoce des TSA que pour approfondir la compréhension de ce trouble, en offrant la possibilité de repérer des profils neurofonctionnels équivalents.
Dans cette perspective, la littérature récente mobilise des approches fondées sur l'apprentissage automatique pour, d'une part, caractériser la motricité fine des individus avec un TSA, et, d'autre part, classifier les profils neurotypiques de ceux des enfants présentant un TSA.

\subsection{Apport des approches par apprentissage automatique}

\subsubsection{Classification en fonction de variables motrices}\label{sec:classification}

Tout d'abord, la littérature montre de manière convergente la faisabilité d'une classification des TSA par des approches d'apprentissage automatique appliquées à des tâches de motricité fine. Les précisions sont généralement comprises entre 70 et 90\,\% selon la tâche et le modèle utilisés, qu'il s'agisse de tâches d'imitation gestuelle au niveau de la main \parencite{Li2017Aug, FulinLiu2024Apr, vabalas_applying_2020}, de saisie manuelle \parencite{crippa_use_2015,Simeoli2021May,Milano2023May,Su2024Dec, freud_effective_2025, cavallo_identifying_2021}, de dessin \parencite{Emanuele2021Aug}, ou de pointage \parencite{Anzulewicz2016Aug, Luongo2024Apr, Doctor2025Jul}.

\subsubsection{Caractérisation de la motricité fine}
Par ailleurs, les modèles issus de ces études mettent en évidence des différences dans les vitesses, accélérations et dans la variabilité des trajectoires et amplitudes, traduisant des mouvements moins fluides chez les enfants avec TSA \parencite{Simeoli2021May, Milano2023May, Li2017Aug, Doctor2025Jul}. Il a été proposé que les enfants avec TSA auraient une moindre dépendance au \textit{feedback} sensoriel, en particulier visuel, et des stratégies de contrôle davantage en boucle ouverte \parencite{Emanuele2021Aug, crippa_use_2015, Anzulewicz2016Aug}, ce qui est cohérent avec les conclusions des études plus observationnelles (voir Sections \ref{sec:pointage}, \ref{sec:ecriture}).

%lles décrivent la motricité des TSA mais de part le fait que les tâches et variables aient été choisie a priori on sait pas si elles sont représentatives, on sait pas si elles sont les plus caractéristiques + on peut pas les comparer entre elles !!

\medskip

Pour autant, les approches par apprentissage statistique demeurent à un stade de preuve de concept \parencite{Doctor2025Jul}. Bien qu'il soit possible d'extraire de l'information pertinente pour la classification à partir de tâches de motricité fine, les variables de confusion, telles que l'âge ou la présence ou non de déficience intellectuelle, peuvent biaiser l'apprentissage des modèles et ne sont pas toujours prises en compte \parencite{cavallo_identifying_2021}. De plus, les échantillons étudiés sont souvent relativement faibles comparés au nombre de \textit{features}, ce qui entraîne un risque de sur-apprentissage%[source,]
. Enfin et surtout, la validation d'une approche automatique reste limitée par une forte hétérogénéité méthodologique, tant au niveau des tâches que des \textit{features} extraites, ce qui empêche la constitution d'un modèle entraîné sur un grand volume de données et donc pleinement généralisable.

Face à cette hétérogénéité, l'identification des variables motrices les plus caractéristiques permettrait d'orienter la collecte de données pour le développement d'outils diagnostics basés sur la motricité. De plus, dans une perspective de modélisation \textit{utile}, cela permettrait de formuler des hypothèses sur les systèmes neuronaux sous-jacents impliqués chez les sujets TSA. Par exemple, si les variables les plus discriminantes reflètent principalement des processus de planification, cela suggérerait une implication prédominante des ganglions de la base et permettrait de mobiliser les connaissances issues de la maladie de Parkinson pour mieux comprendre les mécanismes autistiques \parencite{Rinehart2006Aug}.

% Problématique, objectifs et hypothèses
\chapter{Problématique, objectifs et hypothèses}

\section{Problématique}
La revue de la littérature met en évidence le besoin de mieux comprendre les TSA et de pouvoir les diagnostiquer précocement. Bien que l'on possède des connaissances sur les bases neuronales, la compréhension des TSA et les outils de diagnostic actuels restent largement imparfaits. En ce sens, une caractérisation de la motricité fine des TSA, notamment par des approches computationnelles, pourrait être une stratégie pertinente pour combler ces limites. Toutefois, ces approches ne sont pas cliniquement exploitées. Une limite majeure que l'on peut noter est l'hétérogénéité des variables mesurées et des tâches utilisées. Ainsi, déterminer les marqueurs moteurs les plus caractéristiques des TSA constituerait une avancée pour la compréhension et le diagnostic de ce trouble. 

\bigskip

Nous nous proposons d'identifier, parmi un ensemble de tâches de motricité fine, les marqueurs moteurs les
plus caractéristiques chez l'enfant avec TSA, dans l'objectif de mieux comprendre les altérations motrices associées et de contribuer, à titre exploratoire, au développement d'outils diagnostiques complémentaires à ceux utilisés aujourd'hui.

\section{Objectifs}
%L'objectif principal de ce mémoire est de déterminer les marqueurs moteurs les plus caractéristiques des TSA chez l'enfant, à partir de données issues de tâches de motricité fine. 
À l'aide d'analyses adaptées à un contexte de forte dimensionnalité, ce travail vise, d'une part, à améliorer la compréhension des TSA en reliant les marqueurs moteurs identifiés aux corrélats neuronaux décrits dans la littérature. D'autre part, il cherche à identifier une tâche ou un ensemble de variables de motricité fine présentant un fort pouvoir discriminant, afin d'évaluer l'existence de différents profils moteurs au sein de la population TSA. Cette approche pourrait constituer une base pour de futures études à plus large échelle, visant le développement d'outils de classification automatique en complément du diagnostic actuel.


\section{Hypothèses}
Sur la base de la littérature et des questionnements précédents, nous faisons l'hypothèse qu'il sera possible de classifier les enfants présentant un TSA par rapport aux enfants neurotypiques avec une précision de l'ordre de 70 à 90\,\% -- avec une sensibilité et spécificité équilibrée -- à partir de l'ensemble des variables motrices issues des différentes tâches, et ce, sans sur-apprentissage. Nous supposons par ailleurs qu'un sous-ensemble restreint, correspondant aux variables les plus discriminantes, sera suffisant pour classifier les deux groupes avec des performances comparables à celles obtenues avec l'ensemble complet des marqueurs% (typiquement <1 sd ELPD/qq \%)
. Enfin, nous faisons l'hypothèse que les variables motrices les plus discriminantes présenteront des liens cohérents avec les corrélats neurophysiologiques décrits dans la littérature, principalement le cervelet et les ganglions de la base, et qu'elles seront donc majoritairement associées à la fluidité et à l'accélération du mouvement.


%hypothèses trop générales ?
%partie moteur/neuro trop rapide ?